% GENERATED FILE. Edit canonical lessons, then run npm run generate:books.
% canonical-source-hash: fnv1a64:22c0982b9a30da92
% canonical-lessons: ZH-C223-dongwu, ZH-C223-xiongmao, ZH-C223-laohu, ZH-C223-houzi, ZH-C223-tuzi, ZH-R223-review-218-223

\chapter{Animals}
\label{ch:zh-animals}
\hlchaptermodality{223}

\begin{chapteropening}
\textbf{By the end of this chapter:} \emph{I can name animals: a panda, a tiger, a monkey and a rabbit, and say which ones I like.}

Complete the last lesson of chapter 223: I can name animals: a panda, a tiger, a monkey and a rabbit, and say which ones I like.
\end{chapteropening}

\section[dòngwù]{dòngwù — animal}

\label{lesson:ZH-C223-dongwu}



\begin{quote}
Before the new one: say the Mandarin for to sneeze, then the Mandarin for miserable.
\end{quote}

\subsection*{You'll want to know: dòngwù}
\emph{dòngwù} — \textquotedblleft{}animal\textquotedblright{}.

\begin{quote}
\emph{nǐ xǐhuan shénme dòngwù} — What animals do you like?
\end{quote}

Its characters come later; for now it is a word for the ear.

\begin{cousinweb}
\emph{dòngwù} is built from \emph{dòng} 'to move' and \emph{wù} 'thing', both inherited from Middle Chinese, \emph{duwngX} and \emph{mjut} (Guangyun, in Baxter's transcription), which glosses the first 'to shake' and the second 'the ten thousand things'.
\end{cousinweb}

\begin{grammarlens}[title={What you've built}]
One more word for everyday things.
\end{grammarlens}

\subsection*{Your turn}
\begin{itemize}
\raggedright
  \item \emph{Say it:} \emph{dòngwù}
  \item \emph{Say it:} \emph{dòngwù}, once more
  \item \emph{Say it:} \emph{nǐ xǐhuan shénme dòngwù} — What animals do you like?
  \item \emph{Recall:} say the Mandarin for to sneeze, then the Mandarin for miserable, then say \emph{dòngwù} again
\end{itemize}

\begin{tcolorbox}[breakable,colback=teal!4,colframe=teal!35!black,title={Before you move on}]
What does dòngwù mean? (\textquotedblleft{}Animal\textquotedblright{}.) Say it once more.
\end{tcolorbox}
\section[xióngmāo]{xióngmāo — a panda; giant panda}

\label{lesson:ZH-C223-xiongmao}



\begin{quote}
Before the new one: say the Mandarin for miserable, then the Mandarin for an animal.
\end{quote}

\subsection*{You'll want to know: xióngmāo}
\emph{xióngmāo} — \textquotedblleft{}a panda; giant panda\textquotedblright{}.

\begin{quote}
\emph{xióngmāo hěn kě'ài} — Pandas are so cute.
\end{quote}

Its characters come later; for now it is a word for the ear.

\begin{cousinweb}
\emph{xióngmāo} is built from \emph{xióng} 'bear' and \emph{māo} 'cat', both inherited from Middle Chinese, \emph{hjuwng} and \emph{mjew} (Guangyun, in Baxter's transcription), which lists the first as a beast like a pig and the second as the animal that catches mice; the level tone of māo is not the regular outcome of that form.
\end{cousinweb}

\begin{grammarlens}[title={What you've built}]
One more word for everyday things.
\end{grammarlens}

\subsection*{Your turn}
\begin{itemize}
\raggedright
  \item \emph{Say it:} \emph{xióngmāo}
  \item \emph{Say it:} \emph{xióngmāo}, once more
  \item \emph{Say it:} \emph{xióngmāo hěn kě'ài} — Pandas are so cute.
  \item \emph{Recall:} say the Mandarin for miserable, then the Mandarin for an animal, then say \emph{xióngmāo} again
\end{itemize}

\begin{tcolorbox}[breakable,colback=teal!4,colframe=teal!35!black,title={Before you move on}]
What does xióngmāo mean? (\textquotedblleft{}A panda; giant panda\textquotedblright{}.) Say it once more.
\end{tcolorbox}
\section[lǎohǔ]{lǎohǔ — tiger}

\label{lesson:ZH-C223-laohu}



\begin{quote}
Before the new one: say the Mandarin for an animal, then the Mandarin for a panda.
\end{quote}

\subsection*{You'll want to know: lǎohǔ}
\emph{lǎohǔ} — \textquotedblleft{}tiger\textquotedblright{}.

\begin{quote}
\emph{wǒ mèimei pà lǎohǔ} — My little sister is scared of tigers.
\end{quote}

Its characters come later; for now it is a word for the ear.

Tone: lǎo and hǔ are both third tones, so the first rises and it sounds like láohǔ.

\begin{cousinweb}
\emph{lǎohǔ} is built from \emph{lǎo} 'old' and \emph{hǔ} 'tiger', both inherited from Middle Chinese, \emph{lawX} and \emph{xuX} (Guangyun, in Baxter's transcription), which calls the second, after an older book, the lord of the mountain beasts. Here the first is a prefix with no sense of age; how it came to sit there is not settled.
\end{cousinweb}

\begin{grammarlens}[title={What you've built}]
One more word for everyday things.
\end{grammarlens}

\subsection*{Your turn}
\begin{itemize}
\raggedright
  \item \emph{Say it:} \emph{lǎohǔ}
  \item \emph{Say it:} \emph{lǎohǔ}, once more
  \item \emph{Say it:} \emph{wǒ mèimei pà lǎohǔ} — My little sister is scared of tigers.
  \item \emph{Recall:} say the Mandarin for an animal, then the Mandarin for a panda, then say \emph{lǎohǔ} again
\end{itemize}

\begin{tcolorbox}[breakable,colback=teal!4,colframe=teal!35!black,title={Before you move on}]
What does lǎohǔ mean? (\textquotedblleft{}Tiger\textquotedblright{}.) Say it once more.
\end{tcolorbox}
\section[hóuzi]{hóuzi — monkey}

\label{lesson:ZH-C223-houzi}



\begin{quote}
Before the new one: say the Mandarin for a panda, then the Mandarin for a tiger.
\end{quote}

\subsection*{You'll want to know: hóuzi}
\emph{hóuzi} — \textquotedblleft{}monkey\textquotedblright{}.

\begin{quote}
\emph{nà zhī hóuzi hěn cōngming} — That monkey's very clever.
\end{quote}

Its characters come later; for now it is a word for the ear.

\begin{cousinweb}
\emph{hóuzi} is built from \emph{hóu} 'monkey' and the noun ending \zh{子}, both inherited from Middle Chinese, \emph{huw} and \emph{tsiX} (Guangyun, in Baxter's transcription), which glosses the first with a word for the macaque.
\end{cousinweb}

\begin{grammarlens}[title={What you've built}]
One more word for everyday things.
\end{grammarlens}

\subsection*{Your turn}
\begin{itemize}
\raggedright
  \item \emph{Say it:} \emph{hóuzi}
  \item \emph{Say it:} \emph{hóuzi}, once more
  \item \emph{Say it:} \emph{nà zhī hóuzi hěn cōngming} — That monkey's very clever.
  \item \emph{Recall:} say the Mandarin for a panda, then the Mandarin for a tiger, then say \emph{hóuzi} again
\end{itemize}

\begin{tcolorbox}[breakable,colback=teal!4,colframe=teal!35!black,title={Before you move on}]
What does hóuzi mean? (\textquotedblleft{}Monkey\textquotedblright{}.) Say it once more.
\end{tcolorbox}
\section[tùzi]{tùzi — rabbit; hare}

\label{lesson:ZH-C223-tuzi}



\begin{quote}
Before the new one: say the Mandarin for a tiger, then the Mandarin for a monkey.
\end{quote}

\subsection*{You'll want to know: tùzi}
\emph{tùzi} — \textquotedblleft{}rabbit; hare\textquotedblright{}.

\begin{quote}
\emph{wǒ jiā yǒu liǎng zhī tùzi} — We've got two rabbits at home.
\end{quote}

Its characters come later; for now it is a word for the ear.

\begin{cousinweb}
\emph{tùzi} is built from a word for the hare and the noun ending \zh{子}. The noun ending is inherited from Middle Chinese \emph{tsiX} (Guangyun, in Baxter's transcription), and the dictionary lists the first as the name of a beast.
\end{cousinweb}

\begin{grammarlens}[title={What you've built}]
One more word for everyday things.
\end{grammarlens}

\subsection*{Your turn}
\begin{itemize}
\raggedright
  \item \emph{Say it:} \emph{tùzi}
  \item \emph{Say it:} \emph{tùzi}, once more
  \item \emph{Say it:} \emph{wǒ jiā yǒu liǎng zhī tùzi} — We've got two rabbits at home.
  \item \emph{Recall:} say the Mandarin for a tiger, then the Mandarin for a monkey, then say \emph{tùzi} again
\end{itemize}

\begin{tcolorbox}[breakable,colback=teal!4,colframe=teal!35!black,title={Before you move on}]
What does tùzi mean? (\textquotedblleft{}Rabbit; hare\textquotedblright{}.) Say it once more.
\end{tcolorbox}
\section[(dialogue)]{Looking back: from Wind and Warm Days to Animals}

\label{lesson:ZH-R223-review-218-223}



\begin{quote}
Say the Mandarin for a windbreaker and pleasantly warm.
\end{quote}

\subsection*{Your turn}
\begin{itemize}
\raggedright
  \item \emph{Say it:} a windbreaker, pleasantly warm, hot and stuffy, dry air, damp
  \item \emph{Say it:} a cloud, cloudy, to turn colder, air temperature, to clear up
  \item \emph{Say it:} to work out, to lose weight, to sweat, to skip with a rope, to do a fitness routine
  \item \emph{Say it:} a snowman, to ski, ice, below zero, gloves
  \item \emph{Say it:} to catch a chill, to cough, throat, to sneeze, miserable
  \item \emph{Say it:} an animal, a panda, a tiger, a monkey, a rabbit
  \item \emph{Read it:} \zh{风} \zh{云} \zh{雪} — each one aloud, then a word that holds it
\end{itemize}

\begin{tcolorbox}[breakable,colback=teal!4,colframe=teal!35!black,title={Before you move on}]
A windbreaker? (\textbf{\zh{风衣}}, \emph{fēngyī}.) Pleasantly warm? (\emph{nuǎnhuo}.) Hot and stuffy? (\emph{mēnrè}.) Dry air? (\emph{gānzào}.) Damp? (\emph{cháoshī}.) A cloud? (\textbf{\zh{云}}, \emph{yún}.) Cloudy? (\textbf{\zh{多云}}, \emph{duōyún}.) To turn colder? (\emph{jiàngwēn}.) Air temperature? (\emph{qìwēn}.) To clear up? (\emph{fàngqíng}.) To work out? (\emph{jiànshēn}.) To lose weight? (\emph{jiǎnféi}.) To sweat? (\emph{chūhàn}.) To skip with a rope? (\emph{tiàoshéng}.) To do a fitness routine? (\emph{zuòcāo}.) A snowman? (\textbf{\zh{雪人}}, \emph{xuěrén}.) To ski? (\emph{huáxuě}.) Ice? (\emph{bīng}.) Below zero? (\emph{língxià}.) Gloves? (\emph{shǒutào}.) To catch a chill? (\emph{zháoliáng}.) To cough? (\emph{késou}.) Throat? (\emph{sǎngzi}.) To sneeze? (\emph{dǎ pēntì}.) Miserable? (\emph{nánshòu}.) An animal? (\emph{dòngwù}.) A panda? (\emph{xióngmāo}.) A tiger? (\emph{lǎohǔ}.) A monkey? (\emph{hóuzi}.) A rabbit? (\emph{tùzi}.)
\end{tcolorbox}
